\clearpage
\fancyhead[L]{\small GAUSS TO BODE.}
\hypertarget{native-283-lower}{}
\begin{center}
GAUSS TO BODE.
\end{center}

\begin{center}
Astronomisches Jahrbuch für 1809. Pp. 215–219. Berlin 1806.
\end{center}

\begin{flushright}
Brunswick, 30 July 1806.
\end{flushright}

\par\indent In my last letter I promised Your --- that I would send you the results of my latest calculations on the new planets as soon as I had finished them. I now hasten to keep that promise.

\par\indent I have already reported to you my observations of Pallas from February of this year (see \href{https://www.digitale-sammlungen.de/en/view/bsb10538338?page=147}{p.~137} [\hyperlink{native-272}{p.~267} of this volume]). For the latest improvement to the elements, however, I used only the observations made in Milan, which have now also been printed in the \href{https://www.digitale-sammlungen.de/en/view/bsb10538606?page=100}{July issue of the M.\,C.} These, then, are my latest elliptic elements for Pallas.

\begin{center}
\emph{Elements of Pallas.} (IX.)

Epochs of longitude, Seeberg meridian.
\end{center}

\begin{center}\begingroup
\fontsize{9}{11}\selectfont
\setlength{\tabcolsep}{3pt}\renewcommand{\arraystretch}{1.12}
\begin{tabular}{ll}
1802.  & $143^\circ\,24'\,1''1$ \\
1803.  & $221\;32\;56.4$ \\
1804.  & $299\;54\;42.5$ \\
1805.  & $18\;3\;37.8$ \\
1806.  & $96\;12\;33.2$ \\
1807.  & $174\;21\;28.5$ \\
\end{tabular}\endgroup\end{center}
\begin{center}\begingroup
\fontsize{9}{11}\selectfont
\setlength{\tabcolsep}{3pt}\renewcommand{\arraystretch}{1.12}
\begin{tabular}{ll}
daily mean motion .\ .\ .\ .\ . & $770''7816$ \\
Aphelion 1806 .\ .\ .\ .\ .\ .\ .\ . & $301^\circ\,12'\,21''$ \\
ascending node 1806 .\ .\ .\ . & $172\;30\;56$ \\
Inclination of the orbit .\ .\ .\ .\ .\ .\ . & $34\;37\;8$ \\
Eccentricity .\ .\ .\ .\ .\ .\ .\ .\ .\ . & $0.245384$ \\
Logarithm of the semiaxis .\ .\ .\ .\ . & $0.4421021$ \\
\end{tabular}\endgroup\end{center}
\SourceLeaf{284}%
\par\indent Using these elements, I calculated the course of Pallas for its next apparition. I have the honor of enclosing the ephemeris here. Prof. \textsc{Harding} has again made a chart to accompany it, which Mr. \textsc{v. Zach} will have engraved for the M.\,C. A couple of sheets of Prof. \textsc{Harding}’s atlas, which is so necessary for observing the asteroids, will appear shortly.

\par\indent I have already sent you the observations of Juno made here as well (see \href{https://www.digitale-sammlungen.de/en/view/bsb10538338?page=147}{p.~137} [\hyperlink{native-272}{p.~267} of this volume]). As far as I can see at present, they were the first this year. Here too, to improve the elements, I relied chiefly on \textsc{Oriani}’s observations printed in the \href{https://www.digitale-sammlungen.de/en/view/bsb10538606?page=98}{July issue of the M.\,C.}

\Needspace{10\baselineskip}
\begin{center}
VI. \emph{Elements of Juno.}

Epochs of mean longitude. Seeberg meridian.
\end{center}

\begin{center}\begingroup
\fontsize{9}{11}\selectfont
\setlength{\tabcolsep}{3pt}\renewcommand{\arraystretch}{1.12}
\begin{tabular}{ll}
1805.  & $42^\circ\,35'\,7''3.$ \\
1806.  & $125^\circ\,11'\,20''1.$ \\
1807.  & $207^\circ\,47'\,33''0$ \\
\end{tabular}\endgroup\end{center}

\begin{center}\begingroup
\fontsize{9}{11}\selectfont
\setlength{\tabcolsep}{3pt}\renewcommand{\arraystretch}{1.12}
\begin{tabular}{ll}
Longitude of aphelion in 1806 .\ .\ . & $233^\circ\,17'\,1''$ \\
Longitude of the ascending node in 1806  & $171\;4\;57$ \\
Inclination of the orbit .\ .\ .\ .\ .\ .\ . & $13\;3\;28''4$ \\
daily mean tropical motion . & $814''720$ \\
Eccentricity .\ .\ .\ .\ .\ .\ .\ .\ . & $0.2549441$ \\
Logarithm of the semiaxis .\ .\ .\ . & $0.4260480$ \\
\end{tabular}\endgroup\end{center}

\par\indent I have already sent the ephemeris calculated from these elements to Prof. \textsc{Harding} as well, so that he can draw the chart from it. I enclose a copy for you. It will be extremely difficult, however, to observe Juno next year. It was already extremely faint this year, and next year it will not appear even half as bright as it did this year\HistoricalNote{24}. [The course of Juno from 7 December 1806 to 21 September 1807 is given at three-day intervals for 12 o’clock at Seeberg, with right ascension and declination given to the minute.] The $\text{\AstroSymbols ☍}$ of $\PallasSymbol$ therefore occurs on 4 May 1807 at 21 hours, at longitude $223^\circ\,53'$ and northern latitude $42^\circ\,18'$.

\par\indent I have made no further calculations on Ceres this year. I still lack good observations near opposition. Since $\CeresSymbol$ was in such a favorable position, I have no doubt that this $\text{\AstroSymbols ☍}$ will have been well observed in several places. If you have any such observations, you would do me a favor by sending them to me \emph{soon}. I will then gladly send you the results as well\renewcommand{\thefootnote}{**)}\HistoricalNote{25}\renewcommand{\thefootnote}{*)}.



